% GENEROVANO z ../skripta/preamble.tex a kapitoly/makra.tex -- globalni styly a makra obrazku
% ---------- TikZ styly ----------
\tikzset{
  >={Stealth[round]},
  % paska Turingova stroje
  cell/.style={draw=black!70,minimum width=7mm,minimum height=7mm,inner sep=0pt,
    font=\small\ttfamily,anchor=west,fill=white},
  cellhl/.style={cell,fill=hl},
  cellgr/.style={cell,fill=black!8,text=black!50},
  headarrow/.style={-{Triangle[length=2.2mm,width=2.4mm]},thick,draw=warnc,fill=warnc},
  % automaty / stavove diagramy
  st/.style={state,minimum size=10mm,inner sep=1pt,font=\small,draw=defc!80!black,
    fill=defc!6,thick},
  stacc/.style={st,accepting,draw=acc,fill=acc!10},
  strej/.style={st,draw=rej,fill=rej!8},
  tr/.style={->,thick,draw=black!75},
  lab/.style={font=\scriptsize,fill=white,inner sep=1pt},
  % grafy
  v/.style={circle,draw=black!80,fill=white,thick,minimum size=6.5mm,inner sep=0pt,font=\small},
  vs/.style={circle,draw=black!80,fill=white,thick,minimum size=2.2mm,inner sep=0pt},
  vsel/.style={v,fill=acc!25,draw=acc!80!black,very thick},
  vbad/.style={v,fill=rej!15,draw=rej,very thick},
  e/.style={draw=black!70,thick},
  esel/.style={draw=acc!80!black,line width=2pt},
  ebad/.style={draw=rej,line width=1.6pt,dashed},
  efaint/.style={draw=black!20,thin},
  wt/.style={font=\scriptsize,fill=white,inner sep=1pt},
  % obecne bloky / krabicky
  blok/.style={draw,rounded corners=3pt,thick,align=center,inner sep=5pt,font=\small,fill=white},
  blokd/.style={blok,draw=defc,fill=defc!7},
  bloka/.style={blok,draw=acc,fill=acc!8},
  blokr/.style={blok,draw=rej,fill=rej!7},
  bloki/.style={blok,draw=intc,fill=intc!10},
  flow/.style={->,thick,draw=black!70},
  note/.style={font=\footnotesize\itshape,text=black!65,align=center},
}

\makeatletter
% \tmtape{x}{y}{seznam symbolu}{index hlavy (1..)}  -- vykresli pasku
% symbol "_" se vykresli jako blank; "#" jako krizek
\def\tm@blank{_}\def\tm@sharp{*}
\newcommand{\tmsym}[1]{\edef\tm@tmp{#1}%
  \ifx\tm@tmp\tm@blank\ensuremath{\sqcup}\else\ifx\tm@tmp\tm@sharp\ensuremath{\sharp}\else#1\fi\fi}
\newcommand{\tmtape}[4]{%
  \foreach \s [count=\i] in {#3} {
    \ifnum\i=#4
      \node[cellhl] at ({#1+(\i-1)*0.7},#2) {\tmsym{\s}};
    \else
      \node[cell] at ({#1+(\i-1)*0.7},#2) {\tmsym{\s}};
    \fi
  }
  \ifnum#4>0
    \draw[headarrow] ({#1+(#4-1)*0.7+0.35},{#2-0.85}) -- ({#1+(#4-1)*0.7+0.35},{#2-0.4});
  \fi
}
\makeatother


% Sdilena makra kapitol o NP-uplnosti (prefixy chR = dlazdice/grafy, chZ = gadgety)
% ---------- Makra teto kapitoly (prefix chR) ----------
\definecolor{chRr}{HTML}{FF8080}  % cervena
\definecolor{chRg}{HTML}{D9D9D9}  % seda
\definecolor{chRb}{HTML}{8080FF}  % modra
\definecolor{chRz}{HTML}{80F080}  % zelena
\definecolor{chRo}{HTML}{FFBF80}  % oranzova
\definecolor{chRy}{HTML}{FFFFA0}  % zluta
% Wangova dlazdice: \chRtile{x}{y}{horni}{pravy}{spodni}{levy}{popisek}; (x,y) = levy dolni roh
\newcommand{\chRtile}[7]{%
  \begin{scope}[shift={(#1,#2)}]
    \fill[#3] (0,1)--(1,1)--(0.5,0.5)--cycle;
    \fill[#4] (1,1)--(1,0)--(0.5,0.5)--cycle;
    \fill[#5] (0,0)--(1,0)--(0.5,0.5)--cycle;
    \fill[#6] (0,0)--(0,1)--(0.5,0.5)--cycle;
    \draw[black!80] (0,0) rectangle (1,1) (0,0)--(1,1) (0,1)--(1,0);
    \node[circle,fill=white,inner sep=0.6pt,minimum size=3.6mm,font=\scriptsize\sffamily\bfseries] at (0.5,0.5) {#7};
  \end{scope}}
\newcommand{\chRA}[2]{\chRtile{#1}{#2}{chRr}{chRg}{chRb}{chRz}{A}}
\newcommand{\chRB}[2]{\chRtile{#1}{#2}{chRr}{chRo}{chRr}{chRg}{B}}
\newcommand{\chRC}[2]{\chRtile{#1}{#2}{chRr}{chRy}{chRg}{chRo}{C}}
\newcommand{\chRD}[2]{\chRtile{#1}{#2}{chRb}{chRg}{chRb}{chRz}{D}}
\newcommand{\chRE}[2]{\chRtile{#1}{#2}{chRg}{chRy}{chRg}{chRo}{E}}
\newcommand{\chRF}[2]{\chRtile{#1}{#2}{chRr}{chRo}{chRb}{chRg}{F}}
\newcommand{\chRG}[2]{\chRtile{#1}{#2}{chRg}{chRy}{chRb}{chRo}{G}}
\newcommand{\chRH}[2]{\chRtile{#1}{#2}{chRr}{chRo}{chRg}{chRg}{H}}
% Popisna dlazdice (text v trojuhelnicich): \chRltile{x}{y}{horni}{pravy}{spodni}{levy}{barva H}{barva P}{barva S}{barva L}
\newcommand{\chRltile}[6]{%
  \begin{scope}[shift={(#1,#2)}]
    \draw[black!80,fill=white] (0,0) rectangle (2,2);
    \draw[black!50] (0,0)--(2,2) (0,2)--(2,0);
    \node[font=\scriptsize] at (1,1.6) {#3};
    \node[font=\scriptsize] at (1.6,1) {#4};
    \node[font=\scriptsize] at (1,0.4) {#5};
    \node[font=\scriptsize] at (0.4,1) {#6};
  \end{scope}}
\tikzset{
  chRlit/.style={v,minimum size=8.5mm,font=\scriptsize},
  chRlitsel/.style={chRlit,fill=acc!30,draw=acc!80!black,very thick},
  chRconf/.style={draw=rej!80,dashed,thick},
}

% ---------- lokalni makra a styly kapitoly ----------
\colorlet{chZT}{tT!75}
\colorlet{chZF}{tF!75}
\colorlet{chZB}{tB!70}
\colorlet{chZN}{white}
\colorlet{chZo}{intc!65}
\colorlet{chZc}{cyan!45}
\colorlet{chZp}{thmc!40}
\tikzset{
  chZv/.style={circle,draw=black!80,thick,minimum size=3.4mm,inner sep=0pt},
  chZlv/.style={circle,draw=black!80,thick,minimum size=7mm,inner sep=0pt,font=\small},
  chZside/.style={rounded corners=6pt,draw=none},
  chZag/.style={circle,draw=black,thick,minimum size=4.6mm,inner sep=0pt,font=\scriptsize\bfseries},
  chZset/.style={rounded corners=2pt,opacity=0.35},
}
% male prasatko: barvy l1 l2 l3 a b c d e o (T/F/B/N)
\newcommand{\chZpig}[9]{%
  \node[chZv,fill=chZ#1] (l1) at (0,2.5) {};
  \node[chZv,fill=chZ#2] (l2) at (1.6,2.5) {};
  \node[chZv,fill=chZ#3] (l3) at (2.6,2.5) {};
  \node[chZv,fill=chZ#4] (a) at (0.45,1.85) {};
  \node[chZv,fill=chZ#5] (b) at (1.15,1.85) {};
  \node[chZv,fill=chZ#6] (c) at (0.8,1.3) {};
  \node[chZv,fill=chZ#7] (d) at (1.45,0.7) {};
  \node[chZv,fill=chZ#8] (e) at (2.15,0.7) {};
  \node[chZv,fill=chZ#9] (o) at (1.8,0.15) {};
  \draw[e] (l1)--(a) (l2)--(b) (a)--(b) (a)--(c) (b)--(c) (c)--(d) (l3)--(e) (d)--(e) (d)--(o) (e)--(o);}
% gadget krizeni: 13 barev v poradi C,N,E,S,W,NW,NE,SE,SW,y,x,x',y'
\makeatletter
\newcommand{\chZcross}[2][]{%
  \foreach \c [count=\i] in {#2} {\expandafter\xdef\csname chZcc\romannumeral\i\endcsname{\c}}%
  \begin{scope}[#1]
  \foreach \n/\p/\k in {C/{(0,0)}/i,N/{(0,1)}/ii,E/{(1,0)}/iii,S/{(0,-1)}/iv,W/{(-1,0)}/v,
     NW/{(-1,1)}/vi,NE/{(1,1)}/vii,SE/{(1,-1)}/viii,SW/{(-1,-1)}/ix,
     y/{(0,2.6)}/x,x/{(-2.6,0)}/xi,xp/{(2.6,0)}/xii,yp/{(0,-2.6)}/xiii}
     {\node[chZv,minimum size=4mm,fill=\csname chZcc\k\endcsname] (\n) at \p {};}
  \foreach \u/\w in {C/N,C/E,C/S,C/W,N/E,E/S,S/W,W/N,NW/y,NW/x,NW/N,NE/y,NE/xp,NE/E,
     SE/xp,SE/yp,SE/S,SW/x,SW/yp,SW/W,y/N,x/W,xp/E,yp/S}
     {\draw[e] (\u)--(\w);}
  \end{scope}}
\makeatother

% odkazy uvnitr obrazku se na webu nevyhodnocuji (obrazek se preklada samostatne)
\providecommand{\hyperlink}[2]{#2}
\providecommand{\hyperref}[2][]{#2}
\providecommand{\hypertarget}[2]{#2}
